% Original Prism Review starter for Chemical Reviews, provided under MIT.
% This is a writing draft, not an official Chemical Reviews class or a
% reproduction of the published journal. The installed TeX distribution
% provides achemso and its BibTeX style; no publisher source is copied here.
\documentclass[letterpaper,11pt]{article}
\usepackage[margin=1in]{geometry}
\usepackage{amsmath,amssymb,graphicx,array,booktabs,tabularx,siunitx}
% Package mode does not apply journal-specific class configuration. Set the
% review's bibliography requirements explicitly; achemso supplies natbib.
\usepackage[articletitle=true,doi=true,super=true,biblabel=brackets,
  maxauthors=10,etalmode=truncate]{achemso}
\usepackage[hidelinks]{hyperref}
\providecommand{\doi}[1]{\href{https://doi.org/#1}{\nolinkurl{#1}}}
\sisetup{uncertainty-mode=separate}
\setcounter{secnumdepth}{3}
\setcounter{tocdepth}{2}
\setlength{\emergencystretch}{2em}

\title{Review manuscript title}
\author{First Author\textsuperscript{*} and Second Author\\
  \small Department, Institution, City, Country\\
  \small \textsuperscript{*}Correspondence: \texttt{first.author@example.org}}
\date{}

\begin{document}
\maketitle
\begin{abstract}
Replace this paragraph with a self-contained, single-paragraph English
abstract of 100--200 words. State the subject, scope and objectives of the
review, and identify the major topics it covers. Describe only the content
actually developed in the manuscript. This starter provides a structure for
writing and synthetic demonstrations of figures, tables, equations and
references; it does not summarize a literature search or make a scientific
claim. Establish the review's boundaries and balance before drafting its
summary. Remove these instructions and replace every demonstration before
submission. Follow the current journal guidance and the scope approved by
the editors. An invitation or an approved proposal is required before a
Chemical Reviews manuscript is submitted.
\end{abstract}

\tableofcontents
\clearpage

\section{Introduction}\label{sec:introduction}
Define the purpose, intended readership, scope and limitations of the review,
including the years of literature covered. Explain its relation to earlier
reviews and define any special terminology. Do not claim a literature search,
screening process or coverage period that the authors have not established.
This public citation demonstrates reference handling only; retain it only if
relevant to the actual review\cite{wilkinson2016fair}.

\section{Conceptual foundations}\label{sec:foundations}
\subsection{Terminology and organizing principles}\label{sec:terminology}
Introduce the concepts needed to understand the topic. Replace these example
headings with the review's own logical organization. Explain important
distinctions without assuming that readers already know the specialized field.

\subsection{Scope of the evidence}\label{sec:evidence}
Describe the kinds of published evidence considered and the boundaries of the
comparison. Identify missing or incompatible information accurately; do not
replace missing observations with invented values.

\section{Thematic developments}\label{sec:themes}
\subsection{First theme}\label{sec:first-theme}
Discuss the significant contributions relevant to this theme, including work
from different groups. Distinguish what the cited studies demonstrate from
the review authors' interpretation.

\subsection{Second theme}\label{sec:second-theme}
Organize this section around a useful scientific question or conceptual
connection. Compare the conditions, assumptions and limitations of the cited
work rather than presenting an uncritical sequence of paper summaries.

\subsubsection{A narrower comparison}\label{sec:narrower-comparison}
Use a third-level heading only when it helps readers follow a substantial
discussion. The automatic contents above lists the first two heading levels;
adjust its depth deliberately if a longer review requires more detail.

\section{Critical comparison}\label{sec:comparison}
Synthesize the published evidence and explain agreements, disagreements and
the limits of comparisons. Figure~\ref{fig:comparison} is a two-panel layout
slot. Its boxes and Table~\ref{tab:comparison} contain no scientific finding.

\begin{figure}[htbp]
\centering
\begin{minipage}[t]{0.47\linewidth}
\centering\textbf{a}\par\smallskip
\fbox{\parbox[c][32mm][c]{0.90\linewidth}{\centering
Conceptual comparison\\[5pt]Replace with an original figure}}
\end{minipage}\hfill
\begin{minipage}[t]{0.47\linewidth}
\centering\textbf{b}\par\smallskip
\fbox{\parbox[c][32mm][c]{0.90\linewidth}{\centering
Relationship or framework\\[5pt]Keep panel labels readable}}
\end{minipage}
\caption{\textbf{Synthetic comparison layout.} (a) First conceptual slot.
(b) Related conceptual slot. Replace both panels and this description with
the actual display. Preserve the meaning, labels and units of original
graphics; document attribution and permission when reusing published work.}
\label{fig:comparison}
\end{figure}

\begin{table}[htbp]
\centering\small
\caption{Synthetic values illustrating a comparison table.}
\label{tab:comparison}
\begin{tabularx}{\linewidth}{@{}>{\raggedright\arraybackslash}X
 S[table-format=3.0] S[table-format=1.3(3)] S[table-format=1.3]@{}}
\toprule
Example & {Count} & {Estimate $\pm$ SD} & {Score}\\
\midrule
Example A & 120 & 0.385 +- 0.006 & 0.759\\
Example B & 98 & 0.608 +- 0.051 & 0.206\\
\bottomrule
\end{tabularx}
\par\smallskip\footnotesize\raggedright
All values are invented for typesetting. Replace them with verified
information, identify its sources, and define units and uncertainty.
\end{table}

Equation~\eqref{eq:comparison} demonstrates one numbered multiline expression
only. Adopt an equation only when it is relevant and its symbols, assumptions
and source can be explained accurately.
\begin{equation}
\begin{split}
 s_{\mathrm{A}} &= \sum_{i=1}^{n} w_i x_i,\\
 \Delta s &= s_{\mathrm{A}}-s_{\mathrm{B}}.
\end{split}\label{eq:comparison}
\end{equation}

\section{Challenges and open questions}\label{sec:challenges}
Identify unresolved issues and limitations supported by the literature.
Separate missing evidence, disagreement between studies and questions that
remain speculative. Discuss relevant risks accurately where applicable.

\section{Conclusions and outlook}\label{sec:outlook}
Explain the insights justified by the review and the remaining challenges.
State future directions as questions or proposals rather than completed
experiments or established findings.

\section*{Supporting Information}
% Supply a brief nonsentence contents description, including the file type;
% keep this paragraph before Acknowledgments and References. Remove if absent.
Replace with the actual contents and file type of the supplied material.
The included independent \texttt{supplementary.tex} demonstrates extended
comparison notes and source documentation (PDF).

\section*{Author Information}
\noindent\textbf{Corresponding author:} replace with confirmed contact details.

\subsection*{Author Biographies}
\noindent\textbf{First Author.} Replace with this author's own short biography.
Do not invent positions, qualifications or achievements.

\medskip\noindent\textbf{Second Author.} Replace with this author's own short
biography. Include a biography for every author of the actual review.

\section*{Acknowledgments}
Replace with confirmed acknowledgments and funding, or remove if inapplicable.
\section*{Conflict of Interest}
Replace with the authors' confirmed disclosure.
\section*{Abbreviations}
Replace with an accurate list of abbreviations used in the review, or remove
this section if no list is needed.

% Supply complete, verified English entries, Title Case article titles and
% accurate CASSI journal abbreviations; the style does not create these fields.
\bibliography{references}

\clearpage
\section*{For Table of Contents Only}
% The outside of this placeholder, including padding and rule, is 5 x 5 cm.
% Replace it with the authors' original graphic. The final submission may
% require a separate image file; this starter does not package submission files.
\begin{center}
\fbox{\parbox[c][\dimexpr50mm-2\fboxsep-2\fboxrule\relax][c]
  {\dimexpr50mm-2\fboxsep-2\fboxrule\relax}{\centering
  \sffamily\fontsize{8}{10}\selectfont
  TOC graphic placeholder\par\medskip
  Replace with original artwork\par\medskip
  5 cm $\times$ 5 cm}}
\end{center}
\end{document}
